\section{Mixed Heads Close the Accuracy Gap}
\label{sec:mixed}
%% ============================================================

The structural ablations establish that the $V{=}K$ constraint is the main bottleneck
for perplexity.  A natural fix: replace a fraction of energy heads with standard
GPT heads ($V = W_V x$), sharing a single output projection $W_O$.

\begin{table}[h]
\centering
\caption{Iso-family comparison.  All trained 30k steps.  FLOPs = forward pass, $s=4096$.
  V10/V11 use the mixed-head (additive) block.  V12 is a recurrent GPT baseline.}
\label{tab:mixed_results}
\resizebox{\textwidth}{!}{%
% Auto-generated by experiments/energy-inference/scripts/multi-block-ablation/make_tables.py
% Table: mixed_results
% DO NOT edit by hand — re-run the generator instead.
\begin{tabular}{lcrrrrr}
\toprule
\textbf{Model} & \textbf{E:G} & \textbf{Params} & \textbf{MFLOPs/tok} & \textbf{PPL$\downarrow$} & \textbf{Avg$\uparrow$} & \textbf{GSM8K} \\
\midrule
V0 GPT 12$\times$1 & 0:12 & 162M & 321 & \underline{38.31} & 46.7\% & \textbf{2.20\%} \\
V1 EGPT 12$\times$1 & 0:12 & 143M & 359 & 47.66 & 46.5\% & 1.74\% \\
V10 Mixed 12$\times$1 & 0:12 & 144M & 285 & \textbf{36.99} & \textbf{46.7\%} & 1.52\% \\
\midrule
V11 Mixed 6$\times$2 & 0:12 & 118M & 314 & 40.67 & \underline{46.7\%} & 1.36\% \\
V12 GPT 6$\times$2 & 0:12 & 120M & 321 & 39.98 & 46.4\% & \underline{2.05\%} \\
\bottomrule
\end{tabular}
}
\end{table}

\paragraph{Findings.}
\begin{itemize}
  \item V10 Mixed ties V0 GPT on accuracy (47.2\%) with 12\% fewer FLOPs and 11\% fewer
    parameters, and achieves better PPL (37.0 vs.\ 38.3).
    On a per-FLOP basis, V10 is marginally more efficient than V0 at this scale.
  \item V11 (6$\times$2) ties V12 recurrent GPT,
    confirming that the mixed-head design generalises to recurrent settings.
  \item \textbf{GSM8K regresses.}  Both V10 (1.52\%) and V11 (1.36\%) fall below
    the GPT baseline (2.20\%) on mathematical reasoning.
    The energy $V{=}K$ constraint in the energy heads appears to interfere with
    arithmetic reasoning regardless of the fraction of energy heads.
  \item Higher energy-head fractions (V13 8:4, V14 10:2) perform worse than 6:6 (V10),
    confirming the GPT heads contribute disproportionately to accuracy.
\end{itemize}

\paragraph{A ceiling exists.}
Mixed heads with the $V{=}K$ output rule cannot exceed V0 GPT on accuracy or GSM8K,
despite matching on FLOPs.  The energy heads are contributing \emph{something}
(better PPL, some accuracy) but their output is suboptimal.

%% ============================================================
